% ==================================================================
% Color Charge as Emergent Vertex Identity:
% The hDP Tetrahedral Skeleton and SU(3) as Bookkeeping
% in Conscious Point Physics
% Companion 15 to "Mechanistic Derivation of Relativistic Effects
% via Space Stress Vector (SSV) in the Dipole Sea" (Version 16)
% Version 2 — 19 March 2026
% (Merged from independent Claude and Grok drafts)
% ==================================================================

\documentclass[12pt]{article}

\usepackage{amsmath,amssymb,amsthm}
\usepackage{geometry}
\usepackage{hyperref}
\usepackage{booktabs}
\usepackage{enumitem}
\usepackage{parskip}
\usepackage{array}

\geometry{letterpaper, margin=1in}

\newtheorem{theorem}{Theorem}[section]
\newtheorem{proposition}[theorem]{Proposition}
\newtheorem{corollary}[theorem]{Corollary}
\newtheorem{remark}[theorem]{Remark}
\newtheorem{definition}[theorem]{Definition}

% ── CPP macros ────────────────────────────────────────────────────────────────
\newcommand{\lP}{l_P}
\newcommand{\mP}{m_P}
\newcommand{\PSR}{\mathrm{PSR}}
\newcommand{\SSV}{\mathrm{SSV}}
\newcommand{\LSP}{\mathrm{LSP}}
\newcommand{\eCP}{e\mathrm{CP}}
\newcommand{\qCP}{q\mathrm{CP}}
\newcommand{\eDP}{e\mathrm{DP}}
\newcommand{\qDP}{q\mathrm{DP}}
\newcommand{\hDP}{h\mathrm{DP}}
\newcommand{\ZDC}{\mathrm{ZDC}}
\newcommand{\ZBW}{\mathrm{ZBW}}
\newcommand{\as}{\alpha_s}
\newcommand{\Td}{T_d}
\newcommand{\Hf}{H_4}

\title{\textbf{Color Charge as Emergent Vertex Identity:\\
The hDP Tetrahedral Skeleton, the Shared Origin\\
of Color and Fractional Charge, and SU(3)\\
as Bookkeeping in Conscious Point Physics}\\[6pt]
{\large Companion~15 to ``Mechanistic Derivation of Relativistic Effects\\
via Space Stress Vector (SSV) in the Dipole Sea'' (Version~16)}}

\author{Thomas Lee Abshier, ND\\
Hyperphysics Institute\\
\texttt{https://hyperphysics.com}\\
\texttt{drthomas007@protonmail.com}}

\date{19 March 2026 --- Version~2}

\begin{document}
\maketitle

\noindent\textbf{Keywords:} Conscious Point Physics, color charge,
hDP tetrahedron, quark confinement, baryon structure, vertex identity,
SU(3), Z$_3$ symmetry, color neutrality, gluon octet, fractional
charge, 600-cell, emergent symmetry, DP Sea, charge buffer.

%======================================================================
\begin{abstract}
We derive the observed color charge properties of quarks from the
geometry of the Dipole Sea, without postulating SU(3) as a fundamental
gauge symmetry.  Two distinct but related tetrahedral structures carry
the physical content.

The first is the \emph{individual qCP cage}: a regular tetrahedral
($T_d$) substructure of the 600-cell that every quark Conscious Point
occupies.  This cage already gave fractional electric charges
($e/3$, $2e/3$) via 3D projection (companion~2).  Here we show that
the \emph{same cage} also encodes the three color states via its C$_3$
rotational degeneracy: a 120$^\circ$ rotation around the
apex-to-centroid axis cyclically permutes the three base vertices,
generating the color triplet (Theorem~\ref{thm:color_triplet}).
Fractional electric charge and color charge are therefore two
properties of one geometric object — projective and rotational views
of the same tetrahedron.

The second is the \emph{hDP tetrahedral skeleton}: a spontaneously
nucleated structure of hybrid Dipole Particles in the DP Sea that
provides the charge-buffering scaffold for baryons.  Three quarks
occupy three of its four vertices; the fourth remains open, giving
the baryon its net charge.  Color charge, in this picture, is
\emph{vertex identity}: a quark's color is which of the three
occupied tetrahedral vertices it is bonded to.

The central theorem of the paper is: color charge is not an intrinsic
quantum number; it is a relational property describing a quark's
position in a physical bonding structure.  Color neutrality, the
prohibition on free quarks, and the restriction to specific hadron
configurations ($q\bar{q}$, $qqq$, $\bar{q}\bar{q}\bar{q}$) all
follow from charge geometry alone.  The 1+3+4 = 8 qDP chain layer
modes of companion~14 are identified as the geometric origin of the
8 gluon generators; their explicit mode-to-generator correspondence
is deferred to companion~16.  SU(3) is identified as the permutation
algebra that correctly describes this geometry --- it is correct and
useful, but it is derived from the physical structure rather than
postulated as the foundation of it.

No new postulates beyond those of Version~16 are required.
\end{abstract}

%======================================================================
\section*{Plain Language Summary}
%======================================================================

In standard physics (QCD), quarks carry a property called ``color
charge'' --- red, green, or blue --- and the rules governing color
determine how quarks bind into protons, neutrons, and other hadrons.
Only color-neutral combinations are stable; free quarks are never
observed.  These rules are encoded in the mathematics of SU(3), a
gauge symmetry group that physicists postulate as a fundamental law
of nature.

CPP asks the question one level deeper: \emph{what is the physical
structure that produces these rules?}  The answer involves two
tetrahedral objects.

The first tetrahedron is the cage that every quark sits inside,
built from the geometry of the 600-cell lattice.  This cage already
explained why quarks have fractional electric charges ($1/3$ and
$2/3$ of the electron charge) in companion~2.  It turns out the same
cage also explains color: the three base vertices of the tetrahedron,
which cycle through each other under a 120$^\circ$ rotation, are
the physical basis of the three color states.  Red, green, and blue
are not abstract labels --- they are the names we give to the three
base vertices of a geometric object that quarks are physically
embedded in.

The second tetrahedron is the baryon's bonding skeleton: a
spontaneously formed structure built from hybrid Dipole Particles
($\hDP$s) in the DP Sea.  When three quarks need to coexist in close
proximity inside a proton or neutron, they cannot bond directly
because like-charge quarks repel each other.  The $\hDP$ tetrahedron
solves this problem by providing opposite-charge bonding sites at
each vertex, shielding the quarks from each other's repulsion.  Three
quarks bond to three of the four vertices; the fourth vertex stays
open.  The open vertex is what gives the proton its positive charge.

Color neutrality simply means that all three quark-bonding vertices
are occupied.  Color change (gluon emission and absorption) means a
quark hops from one vertex to another.  SU(3) is the algebra that
correctly counts all the ways this can happen --- but it is a
description of the geometry, not the other way around.

%======================================================================
\section{Introduction}
\label{sec:intro}
%======================================================================

Quantum Chromodynamics (QCD) describes the strong nuclear force
through an SU(3) color gauge symmetry~\cite{gross1973,politzer1973}.
Quarks carry one of three color charges (red, green, blue); gluons
carry color-anticolor combinations; all stable hadrons are
color-neutral.  The structure is mathematically elegant and
experimentally confirmed to high precision.  But SU(3) is a
postulate: it is imposed on the theory from outside, with no
derivation from a more fundamental physical substrate.

Companion~14~\cite{c14} derived quark confinement and the Cornell
potential from the qDP chaining mechanism, without invoking SU(3).
That paper left color charge as an input.  The present companion
closes that gap by identifying the \emph{physical substrate} of
color charge in CPP.

The answer comes in two layers.  The first layer is the individual
quark cage: the same regular tetrahedral substructure of the 600-cell
that already gave fractional electric charges (companion~2~\cite{c2})
now gives the color triplet via its C$_3$ rotational symmetry.
Fractional charge and color charge share a common geometric origin
--- they are projective and rotational properties of the same
tetrahedral object.  This is the content of §\ref{sec:cage_and_color}.

The second layer is the baryon's bonding scaffold: a spontaneously
nucleated $\hDP$ tetrahedral skeleton in the DP Sea that allows
like-charge quarks to coexist by mediating their repulsion.  Color
charge, in this framework, is vertex identity --- the label
identifying which bonding site of this scaffold a quark occupies.
This is the content of §\ref{sec:hdp}--\ref{sec:color_rules}.

Together these two layers produce all QCD color phenomenology without
postulating SU(3).  The algebra of SU(3) describes the resulting
symmetry correctly, but it is a mathematical consequence of the
geometry, not a foundational input.

The present paper establishes color as vertex identity and derives
the color selection rules (§\ref{sec:color_rules}).  It identifies
the 8 gluon modes and their structural origin (§\ref{sec:gluons})
but defers the full mode-to-generator correspondence to
companion~16~\cite{c16}, where the ZBW oscillation modes of the
$\hDP$ skeleton are worked out in detail.

%======================================================================
\section{Two Tetrahedra: Individual Cage and Baryon Skeleton}
\label{sec:two_tets}
%======================================================================

It is essential to distinguish two tetrahedral structures that appear
in the CPP strong-force sector.  They are related but not identical,
and conflating them leads to errors in the physics.

\subsection{The Individual qCP Cage}

Every quark Conscious Point ($\qCP$) sits inside a geometric cage
whose structure is determined by the quark's mass generation.
The lightest quarks (up, down) have minimal or no cage.  The strange
quark has a tetrahedral cage ($\Td$ symmetry, 4 vertices), the charm
quark an icosahedral cage (12 vertices), the bottom quark a
dodecahedral cage (20 vertices), and the top quark a fullerene-like
cage of approximately 60 vertices.  These cages are built from the
600-cell GP shells (Version~16~\cite{v16}) and their ZBW oscillation
energy is what gives each quark generation its characteristic mass
(companion~4~\cite{c4}).

For the lightest quarks (up and down) the individual cage is minimal
and its color and charge geometry is carried primarily by the central
$\qCP$ and its immediate DP cloud.  For heavier quarks the cage
provides additional geometric structure.  The tetrahedral cage of
the strange quark is the clearest example of the individual-cage
C$_3$ symmetry discussed in §\ref{sec:cage_and_color}.

\subsection{The hDP Baryon Skeleton}

Distinct from the individual quark cage is the \emph{hDP tetrahedral
skeleton}: a four-vertex structure spontaneously nucleated from hybrid
Dipole Particles ($\hDP$s) in the DP Sea, around which the three
quarks of a baryon bond.  This skeleton is a property of the
\emph{baryon as a whole}, not of any individual quark.  It is the
structure that allows like-charge quarks (such as the two up quarks
in a proton) to coexist at sub-fm separations without dissociating
under mutual repulsion.

These two tetrahedra are related --- both are embedded in the 600-cell
GP network and both carry the $\Td$ point symmetry --- but they play
different physical roles.  The individual qCP cage determines the
quark's mass and the geometric origin of its fractional charge and
color state.  The hDP baryon skeleton determines the bonding structure
of the hadron and the mechanism by which color neutrality is enforced
physically.

%======================================================================
\section{The Individual Cage: Shared Origin of Charge and Color}
\label{sec:cage_and_color}
%======================================================================

\subsection{Fractional Electric Charge from Cage Projection}

The tetrahedral qCP cage has four vertices.  When the 4D cage
geometry is projected onto 3D space (companion~2~\cite{c2}), the
effective charge that reaches a distant observer is reduced by a
geometric factor determined by the projection of the 4D cage onto
the 3D observation direction.  This projection gives fractional
charges $e_q \in \{e/3,\, 2e/3\}$ for quarks, vs.\ $e$ for the
lepton $\eCP$ which has no cage projection reduction.  The thirds
arise directly from the three-fold geometry of the tetrahedral base
triangle.

\subsection{Color Charge from Cage Rotation}

The same tetrahedral cage that gives fractional charge via projection
also encodes the three color states via rotation.

\begin{theorem}[Color triplet from C$_3$ rotation]
\label{thm:color_triplet}
The three base vertices of the tetrahedral qCP cage are related by
the C$_3$ rotation axis passing through the apex vertex and the
centroid of the base triangle.  A 120$^\circ$ rotation around this
axis cyclically permutes the three base vertices.  These three
vertices, and the cyclic permutation relating them, are the geometric
origin of the color triplet.
\end{theorem}

\begin{proof}
The regular tetrahedron has $T_d$ point symmetry (order 24).  Among
its symmetry elements is a proper 3-fold rotation axis C$_3$ through
each vertex and the centroid of the opposite face.  Consider the axis
through the apex vertex $A$ and the centroid of the base triangle
$(V_1, V_2, V_3)$.  Under the C$_3$ rotation by 120$^\circ$:
\begin{equation}
C_3:\quad V_1 \to V_2 \to V_3 \to V_1.
\label{eq:C3}
\end{equation}
The three base vertices are therefore related by a symmetry of the
cage, but they are not identical --- each has a distinct spatial
position and a distinct relation to the apex.  The three
distinguishable-but-C$_3$-related positions are the geometric
realization of the three color states.  Assigning color labels
red $\equiv V_1$, green $\equiv V_2$, blue $\equiv V_3$, the color
rotation red$\to$green$\to$blue$\to$red is exactly the C$_3$
permutation of Eq.~\eqref{eq:C3}.
\end{proof}

The \emph{color state} of a quark is defined by which base vertex
carries the active qDP chain attachment point.  Rotating the chain
attachment point by 120$^\circ$ changes the color state:
red $\to$ green $\to$ blue $\to$ red.  Color change is a physical
rotation of the chain attachment around the C$_3$ axis of the
tetrahedral cage.

\subsection{Z$_3$ as the Color Quantum Number}

The C$_3$ rotation generates the cyclic group $\mathbb{Z}_3$.  We
assign color charge as a $\mathbb{Z}_3$ quantum number to each base
vertex:
\begin{align}
V_1\,(\text{red})   &: \text{color charge } c = 0, \notag\\
V_2\,(\text{green}) &: \text{color charge } c = 1, \label{eq:Z3}\\
V_3\,(\text{blue})  &: \text{color charge } c = 2. \notag
\end{align}
A color rotation (C$_3$ lattice operation) cycles these labels:
$c \to c + 1 \pmod{3}$.  The center of the SU(3) color group is
exactly $\mathbb{Z}_3$ --- this is not a postulate in CPP, it is the
discrete rotational degeneracy of the tetrahedral cage.  SU(3) is
the simplest Lie group with center $\mathbb{Z}_3$ and a faithful
3-dimensional representation, and the cage geometry selects it
automatically.

\begin{remark}[Shared origin: one cage, two properties]
\label{rem:shared_origin}
Fractional electric charge and color charge are two properties of
the same geometric object --- the tetrahedral qCP cage.  Electric
charge in thirds arises from the projective view: the 4D cage
projected onto 3D space reduces the apparent charge by the
three-fold base-triangle factor.  Color charge arises from the
rotational view: the C$_3$ symmetry axis of the same base triangle
generates the color triplet.  Two distinct measurements of one
tetrahedron give two apparently unrelated quantum numbers.  In CPP
they share a single geometric origin.

This connection is absent from the Standard Model: QCD color and
QED fractional charge are treated as independent quantum numbers
assigned by separate symmetry groups (SU(3) and U(1)).  CPP
unifies them in the cage geometry.
\end{remark}

%======================================================================
\section{The hDP Tetrahedral Skeleton}
\label{sec:hdp}
%======================================================================

\subsection{Hybrid Dipole Particles}

A hybrid Dipole Particle ($\hDP$) is a DP whose two poles combine
electromagnetic ($\eCP$) and strong-force ($\qCP$) character.
Unlike the purely electromagnetic $\eDP$ (charge $\pm e$ at its
poles) or the purely strong $\qDP$ (charge $\pm e_q$), the $\hDP$
presents a mixed charge at each pole: one pole with dominant $\eCP$
character, the other with dominant $\qCP$ character.  This hybrid
nature is what allows the $\hDP$ to serve as a charge buffer between
quarks: it couples to both the electromagnetic SSV of the quark's
electric charge and the strong SSV of the quark's color charge,
mediating both interactions simultaneously.

\subsection{Spontaneous Nucleation}

The Dipole Sea contains $\eDP$, $\qDP$, and $\hDP$ populations.
The $\hDP$ subpopulation spontaneously nucleates tetrahedral
structures via three cooperative mechanisms:

\begin{enumerate}
\item \textbf{Alternating charge geometry (primary stabilizer).}
  A tetrahedron with two $+\hDP$ and two $-\hDP$ vertices in an
  alternating pattern minimizes electrostatic energy relative to
  any other four-vertex charge configuration.  The four vertices
  of a regular tetrahedron sit at mutual distances that balance
  nearest-neighbor attraction and next-nearest-neighbor repulsion
  into a stable potential minimum.  This is the dominant
  stabilizing mechanism --- the $\hDP$ tetrahedron is, first and
  foremost, an electrostatically self-trapped structure.

\item \textbf{ZBW resonance locking (size setter).}
  Each $\hDP$ oscillates at the Zitterbewegung frequency
  (companion~4~\cite{c4}).  When four $\hDP$s are within one
  lattice shell, their ZBW oscillations can phase-lock into a
  collective tetrahedral breathing mode.  This resonance adds
  binding energy and sets the equilibrium size of the tetrahedron
  --- the size at which the ZBW phase-locking is optimal.

\item \textbf{600-cell GP site availability (geometric filter).}
  The four vertices of the $\hDP$ tetrahedron must correspond to
  Grid Points on the 600-cell lattice (Version~16~\cite{v16}).
  The tetrahedral subgroup $T_d \subset H_4$ provides exactly the
  GP configurations that permit tetrahedral formation.  The lattice
  does not drive tetrahedron formation, but it constrains which
  tetrahedra can form, enforcing discrete size quantization.
\end{enumerate}

\begin{remark}[Intermediate stability and suitability]
The $\hDP$ tetrahedron is more stable than a pure $\eDP$ tetrahedron
(which lacks the strong-force binding component) and less stable than
a pure $\qDP$ tetrahedron (which lacks the electromagnetic mediation
component).  This intermediate stability is precisely the property
that makes the $\hDP$ tetrahedron a suitable baryon scaffold: stable
enough to persist on the timescale of hadron dynamics, yet flexible
enough to accommodate the quark bonding transitions that correspond
to color change (gluon emission/absorption).
\end{remark}

\begin{remark}[Thermodynamic status]
The $\hDP$ tetrahedron population exists in one of two regimes: 
dynamic thermal equilibrium at the current ambient temperature 
(a small but nonzero steady-state population maintained by thermal 
fluctuations), or a frozen-out relic population nucleated during 
cosmic cooling when the universe passed through $T \sim \Lambda_{\rm QCD}$.  
Both framings are physically consistent; distinguishing between them 
requires a CPP treatment of the early-universe DP Sea thermodynamics, 
which is an open problem (§\ref{sec:open}).
\end{remark}

\subsection{Geometry of the Baryon Skeleton}

The $\hDP$ tetrahedron has four vertices, six edges, and four faces,
with $T_d$ point symmetry (order 24).  Two of its four vertices carry
net $+\hDP$ pole character; two carry net $-\hDP$ pole character,
in an alternating arrangement.  This alternation is the source of
the charge buffering: a $+\qCP$ quark is attracted to a $-\hDP$
vertex, neutralizing the vertex charge locally while allowing the
quark to sit in stable proximity to other $+\qCP$ quarks at adjacent
vertices.

The three quark-bonding vertices are the three base vertices related
by the C$_3$ axis of the baryon skeleton (distinct from, but
structurally analogous to, the C$_3$ axis of the individual qCP
cage).  The apex vertex of the baryon skeleton is the one that
remains unoccupied --- the source of the baryon's net charge.

%======================================================================
\section{Why Baryons Require the Skeleton; Why Mesons Do Not}
\label{sec:buffer}
%======================================================================

\subsection{The Charge Buffer Theorem}

\begin{theorem}[Baryons require the hDP tetrahedral skeleton]
\label{thm:buffer}
Three quarks that include like-sign charge components cannot form
a stable bound state without the $\hDP$ tetrahedral skeleton as a
charge buffer.
\end{theorem}

\begin{proof}
A proton contains two up quarks ($+2e/3$ each) and one down quark
($-e/3$).  Consider the two up quarks in the absence of the $\hDP$
skeleton.  Their mutual repulsive SSV force at separation $r$ is:
\begin{equation}
F_{\rm rep}(r) = \frac{\as\,\hbar c}{r^2} > 0 \quad
\text{(same-sign quarks, repulsive)}.
\label{eq:rep}
\end{equation}
For the two up quarks to remain bound, a restoring force exceeding
$F_{\rm rep}$ must exist.  The available mechanisms without the
skeleton are:
\begin{itemize}[noitemsep]
\item Direct qDP chain attraction requires opposite-sign quark
  endpoints.  Two $+2e/3$ quarks have no such direct chain.
\item The one down quark ($-e/3$) provides partial attraction to
  each up quark, but its charge magnitude $e/3$ is insufficient
  to overcome the $+2e/3 \times +2e/3$ repulsion at sub-fm
  separations.
\end{itemize}
Therefore, without a mediating structure, the two up quarks
dissociate on a timescale $\tau \sim \hbar/(F_{\rm rep} \cdot l_P)$.
The proton cannot form.

With the $\hDP$ skeleton, each up quark bonds to a $-\hDP$ vertex.
The quark-vertex attraction energy $E_{\rm bond} \sim \as\hbar c/r_0$
exceeds the inter-quark repulsion at the tetrahedral vertex separation
$r_0 \approx 0.26$~fm.  The three-quark configuration is stabilized.
\end{proof}

\begin{remark}[The hDP skeleton as a CPP analog of gluon exchange]
In QCD, the binding of like-charge quarks inside the baryon is
mediated by gluon exchange --- specifically by the off-diagonal
color generators that exchange color charge between quarks.  In CPP,
the $\hDP$ tetrahedral skeleton plays the equivalent role: it
provides opposite-sign bonding vertices that allow like-charge quarks
to remain proximate.  The skeleton is the \emph{physical object}
that the SU(3) gluon algebra is bookkeeping.
\end{remark}

\subsection{Why Mesons Do Not Require the Skeleton}

\begin{proposition}[Mesons bind without the hDP skeleton]
\label{prop:meson_no_skeleton}
A quark-antiquark pair does not require the $\hDP$ tetrahedral
skeleton for binding.
\end{proposition}

\begin{proof}
A quark ($+\qCP$) and antiquark ($-\qCP$) carry opposite charges.
They attract directly via the qDP chain mechanism (companion~14~\cite{c14}),
producing the Cornell potential:
\begin{equation}
V_{\rm meson}(r) = -\frac{\as\,\hbar c}{r} + \sigma\,r.
\end{equation}
No charge buffering is required because there is no like-charge
repulsion to buffer.
\end{proof}

\begin{remark}[Why a meson resists tetrahedral skeleton formation]
If an $\hDP$ tetrahedral skeleton were drawn into the vicinity of a
$q\bar{q}$ meson, its three open quark-bonding vertices would attract
additional quarks from the DP Sea.  These incoming quarks would
convert the two-body meson into a multi-body baryon-like configuration,
breaking the meson's two-body binding structure and releasing energy.
The meson therefore actively resists tetrahedral skeleton incorporation
as a matter of energetics.  The skeleton forms around baryons precisely
because the like-charge repulsion between the baryon's quarks provides
the energy surplus that makes skeleton formation thermodynamically
favorable.
\end{remark}

%======================================================================
\section{The 3-of-4 Occupation Rule and Net Baryon Charge}
\label{sec:occupation}
%======================================================================

\subsection{Why Exactly Three Vertices Are Occupied}

\begin{proposition}[3-of-4 occupation rule]
\label{prop:three_of_four}
In a stable baryon, exactly three of the four $\hDP$ tetrahedral
vertices are occupied by quarks.  Full four-vertex occupation is
thermally unstable.
\end{proposition}

\begin{proof}[Mechanistic argument]
Consider what happens when a fourth quark approaches the unoccupied
fourth vertex of an already three-quark-occupied skeleton.

\textbf{Symmetry equalization.}  When all four vertices are occupied,
the full $T_d$ symmetry of the tetrahedron is restored.  This equal
treatment of all four vertices distributes the available binding
energy uniformly, making each of the four bonds weaker than the
three bonds in the 3-of-4 case.  There is no thermodynamic energy
minimum that uniquely selects four-vertex occupation over three-vertex
occupation plus a free fourth quark.

\textbf{SSV invasion of the fourth vertex.}  The two like-charge
quarks already occupying adjacent vertices (e.g., the two up quarks
at $V_1$ and $V_2$ in a proton) emit repulsive SSV fields.  At the
fourth vertex $V_4$, these fields partially cancel the attractive
SSV of the $\hDP$ vertex charge:
\begin{equation}
E_{\rm bind}(V_4) = E_{\rm vertex}^{\rm attract}
- \sum_{i \in \{\rm like\text{-}charge\}} \frac{\as\,\hbar c}{|V_i - V_4|}
< E_{\rm bind}(V_{1,2,3}).
\label{eq:V4_bind}
\end{equation}
When $E_{\rm bind}(V_4)$ falls below the thermal disruption threshold
$k_B T_{\rm QCD} \approx \Lambda_{\rm QCD} \approx 213$~MeV, the
fourth quark cannot maintain its bond against thermal collisions.

The combined effect of symmetry equalization and SSV invasion ensures
that $E_{\rm bind}(V_4) < k_B T_{\rm QCD}$ for all stable baryon
configurations, leaving the fourth vertex persistently unoccupied.
A rigorous quantitative proof requires explicit computation of
Eq.~\eqref{eq:V4_bind} for each baryon species; this is Open
Problem~\ref{op:fourth_vertex}.
\end{proof}

\subsection{Net Baryon Charge from the Open Vertex}

\begin{proposition}[Net baryon charge]
\label{prop:baryon_charge}
The net charge of a baryon is determined by the unshielded charge
of the unoccupied fourth vertex of the $\hDP$ tetrahedral skeleton.
\end{proposition}

\begin{proof}
At each occupied vertex, the bonded quark presents an opposite-sign
CP that locally neutralizes the $\hDP$ vertex charge.  The
long-range electromagnetic SSV field from an occupied vertex is
therefore zero --- the quark and the vertex together form a locally
neutral pair.  The unoccupied vertex has no such partner and exposes
its full $\hDP$ pole charge to the far field.

\textbf{Proton ($uud$):}
\begin{itemize}[noitemsep]
\item Up quark 1 ($+2e/3$) bonds to a $-\hDP$ vertex: local charge
  $+2e/3 + (-2e/3) = 0$.
\item Up quark 2 ($+2e/3$) bonds to a $-\hDP$ vertex: local charge $0$.
\item Down quark ($-e/3$) bonds to a $+\hDP$ vertex: local charge
  $-e/3 + (+e/3) = 0$.
\item Fourth vertex (unoccupied $+\hDP$ pole): exposed charge $+e/3$.
\end{itemize}
Net proton charge $= (+2e/3) + (+2e/3) + (-e/3) = +e$, confirmed by
the quark sum and consistent with the exposed $+e/3$ vertex charge
contributing to the total~\cite{pdg2022}.

\textbf{Neutron ($udd$):}
\begin{itemize}[noitemsep]
\item Up quark ($+2e/3$) bonds to a $-\hDP$ vertex: local charge $0$.
\item Down quark 1 ($-e/3$) bonds to a $+\hDP$ vertex: local charge $0$.
\item Down quark 2 ($-e/3$) bonds to a $+\hDP$ vertex: local charge $0$.
\item Fourth vertex (unoccupied $-\hDP$ pole): exposed charge $-e/3$.
\end{itemize}
Net neutron charge $= (+2e/3) + (-e/3) + (-e/3) = 0$, with the
$-e/3$ exposed vertex exactly cancelled by the quark-charge
asymmetry~\cite{pdg2022}.
\end{proof}

\begin{remark}[Baryon charge without independent assumption]
The proton charge $+e$ and neutron charge $0$ are derived here from
the charge geometry of the $\hDP$ skeleton without any separate
assumption about ``baryon number'' or ``hypercharge.''  The net
charge follows mechanically from which $\hDP$ vertex is left open
and what quark combination occupies the other three.
\end{remark}

%======================================================================
\section{Color Charge as Vertex Identity}
\label{sec:color_identity}
%======================================================================

\subsection{Definition}

\begin{definition}[Color charge as vertex label]
\label{def:color}
The color charge of a quark in a baryon is its vertex identity:
the label (red, green, or blue) corresponding to the specific
$\hDP$ tetrahedral vertex to which it is bonded.  Color is not an
intrinsic quantum number carried by the quark in isolation; it is
a \emph{relational property} describing the quark's bonding position
within the tetrahedral skeleton.
\end{definition}

The vertex labels correspond to the three $C_3$-related base vertices
of the baryon skeleton, in direct analogy with the C$_3$ axis of the
individual qCP cage established in §\ref{sec:cage_and_color}:
\begin{equation}
C_3: \quad \text{red}\,(c=0) \to \text{green}\,(c=1) \to
\text{blue}\,(c=2) \to \text{red}\,(c=0).
\label{eq:color_cycle}
\end{equation}
A quark's color state can change only by a physical transition ---
detaching from one vertex and attaching to another.  This is a
concrete mechanical event mediated by a $\qDP$ chain oscillation
(a gluon), not an abstract gauge transformation.

\subsection{Color Change as Vertex Transition}

When a quark changes color --- e.g., from red to green --- it
physically moves its qDP chain attachment point from vertex $V_1$
to vertex $V_2$ on the $\hDP$ tetrahedral skeleton.  In QCD language,
this corresponds to emission and absorption of a gluon carrying
color charge $r\bar{g}$ (red-antigreen):
\begin{equation}
q_{\rm red} + g_{r\bar{g}} \;\to\; q_{\rm green}
\quad\Longleftrightarrow\quad
q_{V_1} + \text{(qDP chain mode } V_1 \to V_2\text{)}
\;\to\; q_{V_2}.
\label{eq:color_change}
\end{equation}
The gluon is not a particle carrying abstract color quantum numbers;
it is a $\qDP$ chain oscillation mode that enables a specific vertex
transition.  The full correspondence between specific chain modes
and specific Gell-Mann generators is established in
companion~16~\cite{c16}.

%======================================================================
\section{Color Selection Rules from Vertex Geometry}
\label{sec:color_rules}
%======================================================================

\begin{theorem}[Color neutrality]
\label{thm:color_neutral}
A color-neutral hadron is a configuration in which the sum of
vertex color charges is zero modulo 3, corresponding physically
to either complete three-vertex occupation of the $\hDP$ skeleton
(baryon) or a quark-antiquark direct chain bond (meson).
\end{theorem}

\begin{proof}
By Definition~\ref{def:color} and Eq.~\eqref{eq:Z3}, the three
color charges are 0, 1, 2 in $\mathbb{Z}_3$.

\textbf{Baryon} ($qqq$): three quarks of \emph{distinct} color labels
(red $+$ green $+$ blue $= 0 + 1 + 2 = 3 \equiv 0 \pmod{3}$).
All three quark-bonding vertices are occupied; all vertex charges are
shielded.  The long-range SSV field is zero.

\textbf{Meson} ($q\bar{q}$): a quark (color $c$) and antiquark
(anticolor $\bar{c} = 3-c$ in $\mathbb{Z}_3$): $c + (3-c) = 3
\equiv 0 \pmod{3}$.  No tetrahedral skeleton is needed; the qDP
chain directly connects opposite-sign quarks.

\textbf{Glueball}: pure $\qDP$ loop excitation with no $\qCP$
endpoints.  Net color charge $= 0$ by construction.

In each case, color neutrality is the condition that the total
$\mathbb{Z}_3$ charge is zero.  Configurations with nonzero total
$\mathbb{Z}_3$ charge have unshielded $\hDP$ vertex charges or
open qDP chain endpoints, both of which produce long-range SSV
fields that attract additional quarks until neutrality is restored.
\end{proof}

\begin{remark}[Distinct, not identical colors in a baryon]
The three quarks in a baryon carry \emph{distinct} color labels
(one each of red, green, blue: $c = 0, 1, 2$), summing to
$0 + 1 + 2 = 3 \equiv 0 \pmod 3$.  Three quarks of \emph{identical}
color ($0+0+0 = 0$ or $1+1+1 = 3 \equiv 0$) would also satisfy the
$\mathbb{Z}_3$ condition algebraically, but such configurations are
forbidden by the Pauli exclusion principle applied to the identical
quantum states.  The physical constraint is both the $\mathbb{Z}_3$
neutrality condition \emph{and} the Pauli constraint; the former
selects color-neutral configurations and the latter selects among
them the one with distinct vertex labels.
\end{remark}

\begin{theorem}[Prohibition on free quarks]
\label{thm:no_free_quark}
A single quark cannot exist as a free particle.
\end{theorem}

\begin{proof}
By the confinement theorem of companion~14 (Theorem~1 in~\cite{c14}),
any isolated quark requires a qDP chain extending to infinity, costing
infinite energy.  In the vertex-identity picture, a single quark
bonded to one vertex of a baryon skeleton would expose three
unoccupied vertex charges.  These unshielded charges attract
additional quarks from the DP Sea until three-vertex occupation is
achieved.  Both the chain topology argument (companion~14) and the
vertex-geometry argument (this paper) independently prohibit free
quarks, providing independent and complementary proofs.
\end{proof}

\begin{proposition}[Allowed hadron configurations]
\label{prop:allowed_configs}
The stable color-neutral configurations are exactly: mesons ($q\bar{q}$),
baryons ($qqq$), antibaryons ($\bar{q}\bar{q}\bar{q}$), and glueballs.
\end{proposition}

\begin{proof}[Proof sketch]
Any configuration must satisfy total $\mathbb{Z}_3$ color charge
$= 0 \pmod{3}$ (Theorem~\ref{thm:color_neutral}) and must have no
infinite-energy open chain (companion~14).  The smallest such
configurations are exactly the four listed.  Configurations with
two quarks expose one unshielded vertex and attract a third quark;
four quarks are thermally unstable (Proposition~\ref{prop:three_of_four}).
Exotic hadrons (tetraquarks, pentaquarks) are double-skeleton
configurations possible as transient resonances; their stability
analysis is deferred to Open Problem~\ref{op:exotic}.
\end{proof}

%======================================================================
\section{SU(3) as Emergent Bookkeeping}
\label{sec:su3}
%======================================================================

\begin{proposition}[SU(3) as the algebra of the geometry]
\label{prop:su3_bookkeeping}
The SU(3) color algebra is the mathematical description of the
$\mathbb{Z}_3$ permutation symmetry of 3-of-4 vertex occupation on
the $\hDP$ tetrahedron and the C$_3$ rotational symmetry of the
individual qCP cage.  It is the correct and minimal algebra for this
geometry, but it is derived from the geometry rather than postulated
as its foundation.
\end{proposition}

The derivation proceeds in three steps.

\textbf{Step 1 — Z$_3$ center.}
Both the individual cage C$_3$ rotation (Eq.~\eqref{eq:C3}) and
the baryon skeleton vertex permutation (Eq.~\eqref{eq:color_cycle})
generate $\mathbb{Z}_3$.  The center of SU(3) is exactly
$\mathbb{Z}_3$.  CPP recovers this as the cage's discrete rotational
symmetry rather than postulating it.

\textbf{Step 2 — Triplet and antitriplet representations.}
The three occupied vertices of the baryon skeleton transform as the
fundamental \textbf{3} representation of SU(3) under color rotation.
The three antiquark vertices transform as the conjugate
$\bar{\textbf{3}}$.  Color neutrality requires the tensor product
to contain the singlet \textbf{1}:
\begin{align}
\mathbf{3} \otimes \bar{\mathbf{3}} &= \mathbf{1} \oplus \mathbf{8}
\quad\text{(meson)},\label{eq:meson_decomp}\\
\mathbf{3} \otimes \mathbf{3} \otimes \mathbf{3}
&= \mathbf{1} \oplus \mathbf{8} \oplus \mathbf{8} \oplus \mathbf{10}
\quad\text{(baryon)}.\label{eq:baryon_decomp}
\end{align}
Both decompositions contain the singlet \textbf{1}, confirming that
mesons and baryons are color-neutral in the SU(3) sense --- a
consequence of the geometry, not a separate input.

\textbf{Step 3 — SU(3) as the minimal fitting algebra.}
Any Lie algebra that (a) has center $\mathbb{Z}_3$, (b) has a
faithful 3-dimensional representation, and (c) reproduces the
observed color selection rules must contain SU(3) as a substructure.
SU(3) is the simplest such algebra.  The 600-cell tetrahedral
geometry selects it automatically.

\begin{remark}[What SU(3) explains vs.\ what CPP explains]
SU(3) correctly predicts all color selection rules, hadron multiplet
structure, and the number of gluons (8).  What SU(3) does not explain:
why color charge exists, why there are exactly three colors, why
baryons require color neutrality, and why the color group has center
$\mathbb{Z}_3$ rather than $\mathbb{Z}_2$ or $\mathbb{Z}_4$.  CPP
answers each question mechanically:
\begin{itemize}[noitemsep]
\item Color exists because quarks bond to $\hDP$ tetrahedral vertices.
\item Three colors because the tetrahedron has three quark-bonding
  vertices related by C$_3$ symmetry.
\item Color neutrality because three-vertex occupation shields all
  vertex charges.
\item $\mathbb{Z}_3$ center because the C$_3$ axis of the tetrahedron
  generates the cyclic group of order 3.
\end{itemize}
SU(3) is the right algebra for the geometry.  But the geometry is
more fundamental than the algebra.
\end{remark}

%======================================================================
\section{The Gluon Octet: Structural Identification}
\label{sec:gluons}
%======================================================================

\subsection{Count Correspondence}

The 8-dimensional adjoint representation of SU(3) --- the gluon
octet --- corresponds in CPP to the 8 distinct $\qDP$ chain modes
identified in companion~14~\cite{c14}:
\begin{equation}
\underbrace{1}_{\text{central}} +
\underbrace{3}_{\text{middle}} +
\underbrace{4}_{\text{outer}} = 8 = \dim(\text{SU}(3)\text{ adjoint}).
\label{eq:8count}
\end{equation}
The 8 Gell-Mann generators decompose as 6 off-diagonal (color-changing)
plus 2 diagonal (color-preserving, Cartan subalgebra $T_3$, $Y$).

\subsection{Structural Mode Identification}

Based on the role of each chain layer in the vertex transition
structure of the $\hDP$ skeleton, the following structural
identification is proposed:

\begin{proposition}[Gluon mode structural identification]
\label{prop:gluon_modes}
The 8 qDP chain layer modes map onto the 8 SU(3) generators as follows:
\begin{itemize}
\item \textbf{3 middle-layer chains at 120°, each acting in two
  directions (raising and lowering the color label):}
  These 6 directional modes correspond to the 6 off-diagonal
  Gell-Mann generators ($\lambda_1, \lambda_2, \lambda_4,
  \lambda_5, \lambda_6, \lambda_7$) --- the color-changing
  operators that mediate vertex transitions $V_i \to V_j$.

\item \textbf{Central chain + one geometrically constrained outer
  chain:}
  These 2 modes correspond to the 2 diagonal Cartan generators
  ($\lambda_3 \equiv T_3$, $\lambda_8 \equiv Y$) --- the
  color-preserving operators that distinguish between color states
  without changing them.
\end{itemize}
The remaining 3 outer-layer chains provide the geometric rigidity
that constrains the overall chain topology, ensuring that the 8
physical modes are the complete basis rather than an over-complete
set.
\end{proposition}

\begin{remark}[Why 8 and not 9]
U(3) has 9 generators; SU(3) has 8.  The missing generator
corresponds to the identity transformation --- a uniform phase
rotation that produces no net color change.  In CPP, this would
be a ``breathing mode'' of the entire $\hDP$ tetrahedron in which
all four vertices expand and contract together.  This mode is
suppressed by the 3-of-4 occupation constraint: since the fourth
vertex is always unoccupied, uniform tetrahedral breathing is
geometrically forbidden.  The constraint $\mathrm{det}(U) = 1$
(the SU vs.\ U distinction) is the algebraic expression of the
physical constraint that the fourth vertex remains open.
\end{remark}

\begin{remark}[Deferral to companion~16]
Proposition~\ref{prop:gluon_modes} is a structural identification,
not a rigorous derivation.  Confirming that each middle-layer chain
mode produces exactly the color rotation corresponding to the
associated Gell-Mann generator requires computing the ZBW oscillation
mode decomposition of the $\hDP$ tetrahedral skeleton and mapping
each mode to a specific vertex transition.  This calculation, including
the verification that the resulting commutation relations reproduce
the SU(3) Lie algebra, is the content of companion~16~\cite{c16}.
\end{remark}

%======================================================================
\section{Particle Structure and the Cage Hierarchy}
\label{sec:cages}
%======================================================================

The $\hDP$ tetrahedron of the baryon and the tetrahedral individual
cage of the strange quark are the simplest instances of a broader
pattern in CPP: all massive particles nucleate around successive
geometric cages of DPs on 600-cell shell structures, and mass
increases monotonically with cage complexity.

\begin{table}[ht]
\centering
\caption{Quark, lepton, and boson cage structure in CPP.  The
individual particle cage (column 3) is distinct from the hDP baryon
skeleton (§\ref{sec:hdp}).  Mass increases with cage complexity.}
\label{tab:cages}
\renewcommand{\arraystretch}{1.35}
\begin{tabular}{@{}lllr@{}}
\toprule
Particle & Central CP & Individual cage & Approx.\ mass \\
\midrule
Electron ($e^-$) & $\eCP$ & None (bare) & 0.511~MeV \\
Up quark ($u$) & $\qCP$ & None & $\sim$2.2~MeV \\
Down quark ($d$) & $\qCP$ & Single $\hDP$ radial & $\sim$4.7~MeV \\
Muon ($\mu$) & $\eCP$ & Tetrahedral (4 vertices) & 105.7~MeV \\
Strange quark ($s$) & $\qCP$ & Tetrahedral (4 vertices) & $\sim$95~MeV \\
Tau ($\tau$) & $\eCP$ & Icosahedral (12 vertices) & 1777~MeV \\
Charm quark ($c$) & $\qCP$ & Icosahedral (12 vertices) & $\sim$1275~MeV \\
Bottom quark ($b$) & $\qCP$ & Dodecahedral (20 vertices) & $\sim$4180~MeV \\
Top quark ($t$) & $\qCP$ & Fullerene-like ($\sim$60 vertices) & $\sim$173~GeV \\
\midrule
W boson & None ($\hDP$ loop) & 6-$\hDP$ ribbon/loop & 80.4~GeV \\
Z boson & None & Icosahedral $\hDP$ (12 vertices) & 91.2~GeV \\
Higgs & None & Dodecahedral $\hDP$ (20 vertices) & 125.1~GeV \\
\bottomrule
\end{tabular}
\end{table}

The generational hierarchy emerges directly from the cage progression:
first generation (electron, up/down, photon) --- minimal or no cage;
second generation (muon, strange/charm, W/Z) --- tetrahedral or
icosahedral cage; third generation (tau, bottom/top, Higgs) ---
dodecahedral or fullerene cage.  Each new generation is a particle
that requires an additional geometric shell of DP cage structure to
be stable, and the energy of that cage is the dominant contribution
to the particle's mass.

The bosons (W, Z, Higgs) carry no central unpaired CP and therefore
no intrinsic charge.  They are pure $\hDP$ cage structures whose
geometry determines their mass and interaction cross-section.  The
W boson's six-$\hDP$ ribbon topology gives it its open-ended,
catalyst-like character --- it facilitates charge transfer and flavor
change rather than binding quarks together.

\begin{remark}[Scope of this section]
Table~\ref{tab:cages} is a structural overview establishing the cage
hierarchy context.  The derivation of individual particle masses from
cage ZBW energies, and the formal CPP treatment of the W, Z, and Higgs
interactions, are reserved for dedicated companions in the SM series.
\end{remark}

%======================================================================
\section{Consistency with Previous Companions}
\label{sec:consistency}
%======================================================================

\begin{itemize}[noitemsep]
\item \emph{Fractional charge} (companion~2~\cite{c2}): The shared
  geometric origin of fractional charge (cage projection) and color
  charge (cage C$_3$ rotation) established in §\ref{sec:cage_and_color}
  deepens the result of companion~2.  Color and charge are two views
  of the same tetrahedral cage.

\item \emph{ZBW mass} (companion~4~\cite{c4}): The individual quark
  cage ZBW oscillation energy gives each quark its Compton-scale mass
  contribution, independent of its color/vertex identity.

\item \emph{CP Exclusion Rule} (companion~1~\cite{c1}): The minimum
  PSR ($\geq \lP/2$) governs the scale at which quark-vertex bonds
  form and break, setting the physical scale of the $\hDP$ tetrahedral
  vertex separation.

\item \emph{qDP chaining and Cornell potential} (companion~14~\cite{c14}):
  Color is a layer above the chain mechanism, not in conflict with it.
  The confinement theorem (infinite energy for free quarks) is
  complemented by the vertex-geometry prohibition of
  §\ref{sec:color_rules}.  Both proofs are independent and mutually
  reinforcing.

\item \emph{600-cell GP structure} (Version~16~\cite{v16}): Both the
  individual qCP cages and the $\hDP$ baryon skeleton occupy GP sites
  of the 600-cell.  The $T_d$ tetrahedral subgroup of $H_4$ is the
  common geometric foundation of both.
\end{itemize}

No new postulates beyond those of Version~16 are required.

%======================================================================
\section{Open Problems}
\label{sec:open}
%======================================================================

\begin{enumerate}
\item \label{op:hdp_size}
  \textbf{First-principles derivation of the hDP tetrahedron size.}
  The equilibrium size of the $\hDP$ tetrahedral skeleton is governed
  by the balance between electrostatic attraction and ZBW resonance
  locking.  A quantitative calculation requires the $\hDP$ ZBW
  frequency (from the hybrid CP composition) and the 600-cell GP
  spacing at the relevant shell --- extending the ZBW mass framework
  of companion~4~\cite{c4} to the hybrid DP sector.

\item \textbf{Cosmological status of the hDP tetrahedron population.}
  Whether the baryon skeleton population is in current thermal
  equilibrium or is a cosmological relic frozen out at
  $T \sim \Lambda_{\rm QCD}$ has direct implications for the baryon
  density of the universe.  A CPP treatment requires the full DP Sea
  thermodynamics at strong-force temperatures, connecting to the
  Big Bang baryon asymmetry problem.

\item \label{op:fourth_vertex}
  \textbf{Rigorous proof of the 3-of-4 occupation rule.}
  Proposition~\ref{prop:three_of_four} gives a mechanistic argument.
  The rigorous proof requires an explicit calculation of
  $E_{\rm bind}(V_4)$ in Eq.~\eqref{eq:V4_bind} for each stable
  baryon species, showing the result is below $\Lambda_{\rm QCD}$
  for all cases.

\item \label{op:exotic}
  \textbf{Exotic hadrons from multi-skeleton configurations.}
  Tetraquarks ($qq\bar{q}\bar{q}$) and pentaquarks ($qqqq\bar{q}$)
  have been observed at the LHC~\cite{pdg2022}.  In CPP these
  correspond to double- or multi-$\hDP$-skeleton configurations.
  Their stability conditions, mass spectrum, and decay patterns
  are open problems.

\item \textbf{Gluon mode-to-generator correspondence (deferred to C16).}
  The explicit mapping between the $1+3+4 = 8$ qDP chain layer modes
  and the 8 Gell-Mann matrices, including the derivation of the
  SU(3) commutation relations from chain mode algebra, is the
  content of companion~16~\cite{c16}.

\item \textbf{Embedding of U(1)$_Y \times$ SU(2)$_L \times$ SU(3)
  in $H_4$.}
  The full Standard Model gauge group is the product of three
  factors.  SU(3) is identified in this paper as arising from the
  $T_d$ subgroup of $H_4$.  The SU(2) and U(1) factors should arise
  from other subgroup structures of $H_4$ --- a larger algebraic
  program deferred to the electroweak series (companions~17 and
  beyond).
\end{enumerate}

%======================================================================
\section{Conclusion}
\label{sec:conclusion}
%======================================================================

\begin{enumerate}
\item \textbf{Two distinct tetrahedra} (§\ref{sec:two_tets}): The
  individual qCP cage (sets quark mass and encodes fractional charge
  and color via its geometry) is distinct from the hDP baryon
  skeleton (provides the charge-buffering scaffold for three-quark
  bound states).  Both are tetrahedral; both are embedded in the
  600-cell; they play complementary roles.

\item \textbf{Shared geometric origin of charge and color}
  (Theorem~\ref{thm:color_triplet}, Remark~\ref{rem:shared_origin}):
  Fractional electric charge ($e/3$, $2e/3$) arises from the
  projective view of the tetrahedral qCP cage; color charge arises
  from the rotational view of the same cage.  Two quantum numbers
  from one geometric object.

\item \textbf{Z$_3$ color quantum number} (§\ref{sec:cage_and_color}):
  The C$_3$ rotation of the tetrahedral base vertices generates the
  cyclic group $\mathbb{Z}_3$.  Color charge is a $\mathbb{Z}_3$
  quantum number ($c = 0, 1, 2$) identifying which base vertex the
  quark's chain attachment occupies.

\item \textbf{Charge buffer theorem} (Theorem~\ref{thm:buffer}):
  Like-charge quarks cannot form stable bound states without the
  $\hDP$ tetrahedral skeleton.  Mesons bind without the skeleton
  because opposite-charge quarks attract directly via the qDP chain.

\item \textbf{3-of-4 occupation rule} (Proposition~\ref{prop:three_of_four}):
  Exactly three vertices are occupied in stable baryons.  The fourth
  is destabilized by symmetry equalization and like-charge SSV
  invasion.

\item \textbf{Net baryon charge} (Proposition~\ref{prop:baryon_charge}):
  Proton charge $+e$ and neutron charge $0$ follow from the charge
  of the unoccupied fourth vertex without additional assumption.

\item \textbf{Color as vertex identity} (Definition~\ref{def:color},
  Theorems~\ref{thm:color_neutral}, \ref{thm:no_free_quark}):
  Color is a relational property --- which vertex a quark occupies
  --- not an intrinsic quantum number.  Color neutrality, the
  prohibition on free quarks, and the allowed hadron configurations
  all follow from charge geometry.

\item \textbf{SU(3) as emergent bookkeeping}
  (Proposition~\ref{prop:su3_bookkeeping}): SU(3) is the correct
  and minimal algebra for the $\mathbb{Z}_3$ tetrahedral geometry.
  It is derived from the physical structure rather than postulated
  as its foundation.

\item \textbf{Gluon octet structural identification}
  (Proposition~\ref{prop:gluon_modes}): The 1+3+4 = 8 qDP chain
  layer modes correspond structurally to the 8 Gell-Mann generators
  (6 off-diagonal color-changing + 2 diagonal color-preserving).
  Full derivation deferred to companion~16.
\end{enumerate}

Color charge, in CPP, is not a mystery requiring a new gauge
symmetry.  It is the mechanical consequence of quarks occupying the
vertices of geometric structures that the Dipole Sea spontaneously
provides.  The cage encodes the color triplet in its rotation
symmetry; the skeleton encodes color neutrality in its occupation
rule.  SU(3) describes both correctly --- but the physics comes
first, and the algebra follows.

%======================================================================
\section*{Acknowledgments}
%======================================================================

The conception of the $\hDP$ tetrahedral skeleton as the physical
substrate of baryon structure, the identification of color charge
as an artifact of the resonant bonding modes of quarks on the
tetrahedral skeleton, and the generational cage hierarchy of quarks
and leptons are intellectual contributions of Thomas Lee Abshier, ND,
developed over multiple years of theoretical development.

The formalization of color as vertex identity, the charge buffer
theorem, the 3-of-4 occupation argument, and the net baryon charge
derivation were developed in collaboration with Claude (Anthropic)
on 19~March 2026.  The explicit Z$_3$ color-quantum-number
assignment, the identification of the shared geometric origin of
fractional charge and color from the same tetrahedral cage, and the
gluon mode raising/lowering structural identification were contributed
independently by Grok (xAI) on 19~March 2026.  The merged Version~2
synthesizes both independent analyses.

%======================================================================
\begin{thebibliography}{99}

\bibitem{v16}
T.~L. Abshier,
``Mechanistic Derivation of Relativistic Effects via Space Stress
Vector (SSV) in the Dipole Sea,''
Version~16, 2026 (main paper).

\bibitem{c1}
T.~L. Abshier,
``The Absolute Moment Postulate,''
Version~3, 2026 (companion~1).

\bibitem{c2}
T.~L. Abshier,
``Microscopic Origin of the Dipole Sea Stiffness~$C$,''
Version~2, 2026 (companion~2).

\bibitem{c4}
T.~L. Abshier,
``Inertial Mass from Zitterbewegung,''
Version~2, 2026 (companion~4).

\bibitem{c5}
T.~L. Abshier,
``Newtonian Gravity from Shell-Broadcast SSV Accumulation,''
Version~2, 2026 (companion~5).

\bibitem{c14}
T.~L. Abshier,
``Quark Confinement and the Cornell Potential from qDP Chaining
in Conscious Point Physics,''
Version~1.1, 2026 (companion~14).

\bibitem{c16}
T.~L. Abshier,
``Gluon Modes and the SU(3) Generator Correspondence in CPP,''
2026 (companion~16, in preparation).

\bibitem{gross1973}
D.~J. Gross and F.~Wilczek,
``Ultraviolet Behavior of Non-Abelian Gauge Theories,''
\textit{Phys.\ Rev.\ Lett.}, \textbf{30}, 1343--1346, 1973.

\bibitem{politzer1973}
H.~D. Politzer,
``Reliable Perturbative Results for Strong Interactions?''
\textit{Phys.\ Rev.\ Lett.}, \textbf{30}, 1346--1349, 1973.

\bibitem{pdg2022}
R.~L. Workman \textit{et al.}\ (Particle Data Group),
``Review of Particle Physics,''
\textit{Prog.\ Theor.\ Exp.\ Phys.}, \textbf{2022}, 083C01, 2022.

\end{thebibliography}

\end{document}
